%Capítulo 4

\section{Conclusiones y las recomendaciones}

En este ultimo capitulo se habla sobre las conclusiones obtenidas durante la investigación, su
desarrollo, logros conseguidos y proyectos futuros que puedan surgir de este.

    \textbf{\textit{Descripción}}

En este proyecto se ha diseñado todo un flujo de trabajo que hace uso de la tecnología LoRaWAN, se
ha buscado la generalización par que distintos usuarios, investigadores o empresas puedan usarlo en
sus monitoreos y obtener distintas métricas asi como tomar acciones en base a estos datos. Además de
ello, los programas desarrollados propuestos pueden ser escalables a distintos tipos de
funcionalidades e inclusive tecnologías.

Como primer punto, se hizo una investigación de la tecnología LoRa y LoRaWAN y se resalto el como
esta puede funcionar en equipos mínimos a grandes distancias utilizando pocos recursos. Se procedió
al desarrollo conjunto de una solución embebida al mismo tiempo que se montaba y configuraba un
servidor maestro que puede ser escalable a múltiples servidores con múltiples nodos. Una vez
teniendo todo configurado todo el sistema se pasó a la programación de una capa superior para la
especialización de estos nodos por dos equipos de investigadores distintos demostrando asi su
capacidad de adaptación para futuros proyectos con el minimo esfuerzo.

Teniendo todo el trabajo se pasó al paso de pruebas de todo el flujo de trabajo obteniendo datos
desde múltiples nodos siendo procesados y enviados a los paneles web para finalmente ser desplegados
apropiadamente, demostrando de esta forma la robustez del sistema. Además de la documentación de
todos estos procedimientos no solo en este trabajo si no en un repositorio privado para los
investigadores involucrados en el proyecto.

    \textbf{\textit{Trabajos futuros}}

A pesar de la robustes del flujo de trabajo, hay cabida para múltiples puntos de mejora o distintas
soluciones extras y capas de funcionalidad que pueden ser implementadas en base a distintos
requisitos.
\begin{itemize}
    \item Montaje de un servidor de servidores en la nube.
    \item Orquestación de contenedores asi como balanceo de cargas entre ellos.
    \item Análisis del rendimiento y coste energético del sistema para la búsqueda de optimización en estas áreas. 
    \item Implementar una solución más compacta para continuar simplificando su funcionalidad en usuarios finales.
\end{itemize}
